\section{Exact substring matching without expanding relators}
\label{sec:compressed-matching}

\subsection{What the new stage needs to decide}

Whitehead powers address repeated elementary descent, but not the earlier
backend's inability to compare overlaps once the explicit-letter allowance
is exceeded. The new fallback first handles a precise subproblem: find a
complete occurrence of a donor's recorded spelling, or its inverse, inside
a cyclic target. The substring primitive is complete for its specified
pattern and text. The resulting relator search is restricted: it does not
try every cyclic conjugate of the donor or find the longest partial overlap.

The implemented occurrence table follows the fully compressed pattern-matching
framework of Lifshits~\cite{lifshits-fcpm}. It represents occurrences touching
a grammar cut by a single arithmetic progression, using dynamic programming
on pattern and text rules. The paper's analysis counts operations on string
positions; bit costs require a further factor depending on their binary
length. Independently, Matsubara et al.~\cite{matsubara-compressed-lcs} give a
polynomial algorithm for longest common substrings of two SLP strings.
That stronger primitive would support a complete cyclic-overlap search on
doubled relators, but it is not implemented here. Thus arbitrary compressed
overlap is an available algorithmic direction, not an established exponential
barrier.

\subsection{The implemented occurrence-table invariant}

For a nonempty pattern node $P$ and text node $T=UV$, define
\[
 \mathcal A(P,T)=\{i:0\le i\le |T|-|P|,\quad
                  T[i:i+|P|]=P,\quad i\le |U|\le i+|P|\}.
\]
Positions are zero based, substrings are half open, and touching either
endpoint of the cut is allowed. For a terminal text the cut is zero.
The cut-occurrence lemma states that this set is an arithmetic progression,
including the empty and singleton cases. We store a progression as
$(s,d,c)$: first position, step and cardinality. A singleton has $d=0$.
The table includes empty entries for sufficiently long text nodes; shorter
text nodes are omitted because they cannot contain the pattern.

For a one-letter pattern, a concatenation's only possible occurrences
at its cut are the last letter of its left child and the first letter of
its right child. Both are available from the word arena. A terminal text
requires only one exact letter comparison. The remaining cells are computed
in ascending pattern-node order, so both children of the pattern have
complete rows available. No occurrence list is expanded.

\subsection{An iterative local-window query}

The internal query $\operatorname{Local}(P,T,[\alpha,\beta))$ returns all
occurrences lying entirely within the specified interval, provided that
its width is at most $3|P|$. Bounds are first clipped to the text. The
implementation visits the parse of $T$ using an explicit stack. It obtains
central occurrences from the stored cell and descends into a child only
when the child's intersection with the interval has length at least $|P|$.
Left-child, central and right-child results are emitted in that order.
Occurrences shared between adjacent pieces can only coincide at an endpoint.
The stack avoids dependence on Python's recursion depth.

Write $q=|P|$. Partition occurrence starts into
\[
 i\le\alpha+q\qquad\hbox{and}\qquad i>\alpha+q.
\]
Every occurrence in the first part touches $\alpha+q$, and every occurrence
in the second touches $\alpha+2q$. Indeed, starts lie between $\alpha$ and
$\beta-q\le\alpha+2q$. Each part is therefore one cut-occurrence progression
restricted to an interval of starts. This gives at most two output
progressions without a greedy merge heuristic.

To combine the ordered pieces of one part, the code computes their first
and last positions, the greatest common divisor of their nonzero steps and
offsets from the first position, and their total cardinality after removing
duplicate endpoints. It then checks that this cardinality equals the number
of points in the proposed progression. All supplied points have the same
residue modulo the resulting step. Consequently equality of the cardinalities
certifies that no hole was filled accidentally. A violated invariant raises
an arithmetic error instead of manufacturing a match.

The local traversal has at most a constant number of branches at any level:
its live intervals have disjoint interiors, each has length at least $q$,
and together they fit in an interval of width at most $3q$. Its depth is at
most the number of text rules. Repeated grammar occurrences therefore do
not cause an exponential traversal.

\subsection{Extending a larger pattern child}

Suppose $P=AB$ and $|A|\ge |B|$. If a full occurrence touches text cut $c$,
its $A$ part is contained in
\[
 [c-|P|,\ c+|A|),
\]
a window of width at most $3|A|$. Query that row locally, then restrict
candidate starts so the full occurrence fits in the text and touches $c$.
Each resulting progression describes overlapping copies of the same word
$A$ that share a position. Their union is periodic with the progression's
step, unless it is a singleton.

Let $e$ be the last ending position of those $A$ occurrences. Candidate
endings at most $e-|B|$ leave their entire prospective $B$ block inside
that periodic union. These blocks begin at congruent positions and are
identical. One exact local query therefore either accepts or rejects this
entire first range of candidates. The remaining endings are within $|B|$
of $e$. A local $B$ query on $[e-|B|,e+|B|)$ gives at most two progressions;
intersect them with the candidate endings and shift back by $|A|$.

When the right child is larger, the implementation uses the symmetric
construction on starting positions: prospective left blocks inside the
periodic union are checked together, and the first $|A|$-wide region is
checked locally. Both orientations are exercised by randomized unbalanced
parses. Candidate progression intersections use exact congruences and the
extended-Euclidean modular inverse. Inconsistent residues yield the empty
set. The final pieces describe exactly the full-pattern occurrences that
touch the text cut and are merged with the checked progression operation.

This establishes the cell invariant inductively. To recover the globally
first occurrence, propagate the earliest position of every text node down
the DAG, in reverse topological order. Every nonempty occurrence touches
a cut of some containing parse node (or belongs to a terminal). Taking
the minimum of a node's earliest position plus its first cut-occurrence
position therefore finds the first occurrence in the entire text. Failure
to find any such cell is an exact negative substring answer.

\subsection{Complexity, fast positive queries, and resource semantics}

Let $m,n$ be the numbers of pattern and text rules. There are at most $mn$
cells. Each nontrivial cell uses a constant number of local traversals,
each visiting $O(n)$ parse pieces and performing exact progression arithmetic.
Thus there are $O(mn^2)$ such visits and integer operations, including gcd
and modular-inverse calls. This is a polynomial bit-time algorithm:
positions and counts have $O(n)$ bits when the text has $n$ binary rules,
and gcd and modular inverse are polynomial in that bit length. We do not
identify the Python implementation's charged work counter with a unit-cost
machine model for these big-integer operations. Space is $O(mn+n)$ integer
records, in addition to the input arena. The arena's node allowance also
independently caps the number of live table cells.

Before constructing a full table, the public primitive checks whether the
pattern is a prefix using exact compressed LCP. If not, a bottom-up DAG
pass finds the earliest occurrence of its first signed letter. If there
is no such letter, or insufficient remaining text even at that first
position, absence is proved. Otherwise one exact substring LCP tests that
candidate. Success is necessarily the first match; failure invokes the
complete table. This is an exact fast path with a complete fallback, not
a sample of possible occurrences. Its positive benefit can be large when
the first candidate begins after an exponentially long block.

Table construction, local traversal, grammar operations and the fast path
all share the caller's remaining work and cancellation allowance. Resource
exhaustion raises \code{CompressedLimit}; the group wrapper translates this
to an inconclusive group probe. Partial tables never escape as a negative
substring result. Global cancellation propagates unchanged. The matcher
uses signed letters and exact integer arithmetic, without hashes deciding
word equality or substring occurrence.

\subsection{Integrating whole-donor deletion}

The fallback runs only with relator moves enabled, when current total
expanded length exceeds \code{max\_letters} and the Whitehead test offers
no shortening. This preserves the ordinary compressed descent paths.
Nonempty donors are visited in descending length, with stable slot order
for ties. For each donor, the search tries its spelling and inverse against
distinct target slots at least as long. Per-relator absolute-generator counts
supply a necessary-condition filter: a donor cannot fit in one cyclic turn
of a target with fewer occurrences of any generator.

For donor $D$ and target $R$, searching $D$ in $RR$ finds every complete
occurrence of this particular spelling across a cyclic seam. Since
$|D|\le|R|$, an occurrence found in the second copy has an equivalent start
in the first; reducing its position modulo $|R|$ is valid. The search does
not turn a multiple-wrap match into a legal overlap. If the resulting cyclic
target is $DW$, keeping $D$ and replacing $DW$ by $W$ preserves normal closure:
\[
 \langle\!\langle D,DW\rangle\!\rangle
   =\langle\!\langle D,W\rangle\!\rangle.
\]
Inversion and cyclic conjugation of a relator also preserve normal closure.
The producer checks the entire proposed prefix by an exact LCP before
slicing it away and cyclically reducing the remainder. It records the
existing \code{relator} move, with donor rotation zero and overlap $|D|$.
The independent literal and compressed verifiers already support these
moves and never call the matcher or donor-search helper. A standalone
presentation simplification result is not a knot verdict; production still
reconstructs the full Wirtinger presentation from the original valid PD.

A useful stalled example is
\[
 \langle a,b\mid a^N,\ ba^Nb^{-1}a\rangle,\qquad N=2^h.
\]
For the tested $h=10,100,500$, the previous capped producer cannot proceed.
The new producer finds the donor beginning at target position one. Deletion
leaves a cyclic conjugate of $a$, which eliminates $a$ and empties both
relations. The two-move trace leaves the free generator $b$, as expected
from the presentation itself. No expanded relator is allocated, including
for $N=2^{500}$. This is again an abstract supplied presentation family;
it is not asserted to be a hard knot-diagram family.


\begin{proposition}[A uniform stalled family resolved without expansion]
For every $N\ge2$, the presentation above has no eligible generator
elimination and no shortening Whitehead cut in its initial state. With relator moves enabled and an
explicit overlap allowance below $2N+3$, the preceding internal producer stalls.
Given sufficient compressed work and storage allowances, the new fallback
completes in two moves. For the binary-power grammar at
$N=2^h$, both moves and their discovery use polynomial bit complexity in $h$.
\end{proposition}
\begin{proof}
Neither generator occurs once in either relator. The weighted Whitehead
graph has an edge between $a$ and $a^{-1}$ of weight $2N-1$, and four
weight-one edges joining each of $b,b^{-1}$ to each of $a,a^{-1}$.
An $a$--$a^{-1}$ cut has minimum weight $2N+1$, equal to the degree
of either endpoint: the large edge is mandatory and each of $b,b^{-1}$
contributes at least one more edge. A $b$--$b^{-1}$ cut has minimum weight
two, also equal to the endpoint degree, because the two length-two paths
through $a$ and $a^{-1}$ are edge-disjoint. Thus every optimal Whitehead
change is zero. The old overlap bridge is unavailable above its cap.
The new donor search finds $a^N$ at target position one and leaves the
cyclic word $b^{-1}ab$, which reduces to $a$. One generator elimination
then removes both relators. The initial grammar has $O(h)$ rules. The
first-letter probe finds the exact donor node, and the fixed-rank graph,
compressed slices, cyclic reduction and substitution all operate on a
polynomial-size grammar. No loop iterates over the $N$ letters.
\end{proof}

\subsection{Validation and remaining work}

The maintained tests compare 7340 exhaustive binary pattern/text pairs,
including all text lengths through seven and pattern lengths through four.
Three hundred random signed-letter instances check every stored table cell
against literal cut occurrences and also check local-window results. A
separate 3000-case progression audit checks intersections, including negative
starts and singleton cases. Large tests represent an occurrence progression
with $2^{60}+1$ elements, force a failed first candidate before an exponential
later match, and use an unbalanced 1500-rule grammar. Table, work and global
cancellation limits are exercised. Three hundred fifty random relator tuples
compare the complete recorded-donor search against literal cyclic scans.
Finally, genuine small-unknot traces inflated by a Whitehead power through
$2^{100}$ receive a new donor-deletion move and pass independent compressed
replay; a smaller exponent also passes literal replay.

The implementation removes an expansion barrier for exact substring queries
and enables additional certified search paths. Completing arbitrary cyclic
partial overlaps, reusing occurrence tables across changing presentations,
and bounding the total number of presentation moves remain separate tasks.
In particular, polynomial cost for an occurrence query does not give a
polynomial, quasi-polynomial or general subexponential bound for recognition.

\subsection{Controlled measurements}

The driver is \path{benchmark_compressed_match.py}.
Its historical baseline is commit \code{59054ec761d1};
the full revision is retained in the raw artifact. It measures the
historical search, an identical A/A control and the new search in shuffled
order, with five measured rounds after one excluded warmup. Full queries
include fresh PD validation, recognition and independently reconstructed
certificate replay. There are 225 measured queries on the fifteen-case knot
corpus and thirty additional adaptive group probes on Gordian at letter
caps 4096 and 8192. Every query and probe completes. Full-query certificate
digests agree across the three arms. None of these runs reaches the new
stalled-state fallback, so the following data check preservation of existing
paths and do not demonstrate a knot-corpus speedup.

\begin{center}\small
\begin{tabular}{@{}lrrr@{}}
\toprule
Scope & Before & A/A & Matching \\
\midrule
Small 14, summed medians & 348.682 & 344.734 & 351.222 \\
Gordian, full query & 1976.739 & 1976.758 & 1985.460 \\
Gordian group, cap 4096 & 3682.891 & 3718.799 & 3706.699 \\
Gordian group, cap 8192 & 1146.458 & 1165.369 & 1158.745 \\
\bottomrule
\end{tabular}
\end{center}
Median milliseconds; all knot queries and group probes complete.


Sixty additional measurements construct and simplify the abstract presentation
family above. The old search and its control stall for every tested exponent;
the new search completes in two moves in every run. At $h=500$, median
construction plus search time is 7.712 milliseconds, with 1014 allocated
nonempty grammar nodes and 17325 charged work units. The earliest-letter
probe finds the useful occurrence, so this family does not require the
full occurrence table. These are capacity results on supplied presentations,
not ratios against completed knot recognition.

\begin{center}\small
\begin{tabular}{@{}rlrrr@{}}
\toprule
$h$ in $N=2^h$ & Before/A/A & New (ms) & Moves & Nodes \\
\midrule
8 & stalled & 0.385 & 2 & 30 \\
32 & stalled & 0.862 & 2 & 78 \\
100 & stalled & 1.897 & 2 & 214 \\
500 & stalled & 7.712 & 2 & 1014 \\
\bottomrule
\end{tabular}
\end{center}
Full construction and search on abstract presentations; stalls are incomplete searches, not completed recognition times.


A separate 135-query audit tests the complete matcher. For $N=2^h$ let
$P=(ab)^N$. The first-hit text is $cPc$; the later-hit text is
$ac\,c^N P$, so its first $a$ is a false candidate; the absent text is
$ac(ab)^{N-1}ac$. The expected first positions are respectively $1$, $N+2$
and absent. Each timing includes grammar construction. The full-table arm
bypasses the fast path; the adaptive and A/A arms both use the first-letter
probe with table fallback. All 135 results complete and agree with the
constructed expected answers.

\begin{center}\small
\begin{tabular}{@{}lrrrr@{}}
\toprule
Substring query & $h$ & Full table & Adaptive & A/A \\
\midrule
First candidate matches & 8 & 4.779 & 0.055 & 0.053 \\
Later candidate matches & 8 & 5.267 & 5.290 & 5.278 \\
No occurrence & 8 & 7.111 & 7.067 & 7.089 \\
First candidate matches & 32 & 135.163 & 0.127 & 0.100 \\
Later candidate matches & 32 & 167.487 & 163.197 & 160.450 \\
No occurrence & 32 & 247.027 & 246.467 & 244.871 \\
First candidate matches & 64 & 946.014 & 0.216 & 0.203 \\
Later candidate matches & 64 & 1164.627 & 1193.875 & 1174.301 \\
No occurrence & 64 & 3274.404 & 3301.257 & 3236.010 \\
\bottomrule
\end{tabular}
\end{center}
Median milliseconds including grammar construction. A/A repeats the adaptive query; both positive and negative results are exact.


At $h=64$, the first-hit probe takes about 0.216 milliseconds and avoids
all table cells; direct table construction takes about 946 milliseconds
and stores 2415 cells. The later-hit and absent cases still require full
tables: 4626 and 4556 cells respectively. Their adaptive medians are about
1.194 and 3.301 seconds. Thus the implementation supports exact queries on
exponential strings, but the polynomial fallback has significant constants.
These measurements motivate future table reuse and better string indexes;
they do not justify moving this fallback onto already successful small-word
search paths.

All runs use twenty million charged work units and 100000 arena nodes/table
cells; whole recognition allows 18 seconds, and its group stage 15 seconds.
The explicit overlap allowance for the abstract presentations is four
letters. The raw artifact \code{results/compressed\_match\_20261008.json}
retains the inputs or exact construction, execution order, statuses, timings
and per-query counters. The integrated suite passes 666 tests with Regina
7.4.1 in 140.313 seconds; its log is
\code{synthesis/data/compressed-match-integrated-tests.txt}.

No new delivery was present in \code{docs/incoming} during this iteration.
The suggested \code{openai/math} repository remained at
\code{fd4aeeb2ee4fc729c18d98444fed42fd0529eeeb}. The newly integrated research
is the exact compressed-pattern-matching method cited above; the fuller
longest-common-substring algorithm remains a documented next step.
